% s04-sec3a.tex —— §3 假设、测度与稳定迹公式：3.1 假设、3.2 测度（原文 p.17–25）
\section{假设、测度与稳定迹公式}
\srcpage{17}
\subsection{假设}

设 $G$ 是整体域 $F$ 上的约化群，$F$ 或是 $\mathbb Q$ 的有限扩张，或是有限域 $\kappa$（$|\kappa|=q$）上曲线 $X$ 的有理函数域。记 $\mathbb A_F$ 为 $F$ 的阿代尔环。

自守表示理论分为三个不同的部分：内窥与 $L$-包，它们把一切都归约到拟分裂群上的稳定迹公式；群 $G$ 上的自守表示与其拟分裂内形式上的自守表示的比较，也就是一些定理的证明，其中所谓的 Jacquet–Langlands 定理是最简单的一个，尽管它也属于内窥理论；以及拟分裂群的函子性。既然基本引理已经得到，我们可以期待前两部分的问题迅速进展（[A3]）。本文及后续文章的主要目标，是着手研究函子性——如果上帝允许的话——这将是迄今为止最困难的部分。因此我们假定 $G$ 及其导群 $G_{\mathrm{der}}$ 都是拟分裂的。\footnote{一位审稿人曾对我们在这个假设上的摇摆不定感到不安。理由很简单。从稳定迹公式的终极理论来看，任意群 $G$ 的公式——无论拟分裂与否，只要假定 $G_{\mathrm{der}}$ 单连通——既可以当作 $G$ 的公式来处理，但以其拟分裂形式 $G_{\mathrm{qd}}$ 作为它的一个内窥群；也可以当作 $G_{\mathrm{qd}}$ 的理论来处理，但使用 $G$ 上的函数 ${}^Gf$ 的内窥变换 ${}^{G_{\mathrm{qd}}}f$。因此我们的假设其实是无关紧要的，只是它有时使解释更简单。}

使用 z-扩张——这一概念在 [K1] 中被形式化——可以把问题归约到导群单连通的情形。群 $G$ 的一个 \textbf{z-扩张}，特别地，是一个扩张 $G'\to G$，使得导群 $G'_{\mathrm{der}}$ 单连通，且同态 $G'(F)\to G(F)$ 与 $G'(\mathbb A_F)\to G(\mathbb A)$ 都是满的。在某种意义上，$G$ 上的自守形式理论包含在 $G'$ 的理论中。因此本文假定 $G$ 的导群 $G_{\mathrm{der}}$ 单连通。正特征整体域的情形在 [K1] 中未作处理；由于时间不足，我们在本文中也只能不完整地处理它。

设 $Z$ 为 $G$ 的中心的中性分支。有正合列
\begin{equation}\tag{3.1}
\{1\}\to A\to Z\times G_{\mathrm{der}}\to G\to\{1\},
\qquad A=G_{\mathrm{der}}\cap Z.
\end{equation}

\srcpage{18}
稍后我们将描述 $F$-点在我们所谓的 \textbf{Steinberg–Hitchin 基}中的结构，它其实只是稳定共轭类的集合。Hitchin 的名字与此相连，是因为我们的思考受到 Lie 代数情形的理论影响；但我们所需要的稳定类分类的要点早于这一理论，见 Steinberg 与 Kottwitz 的文章，尽管两人都没有强调线性参数。若有限群 $A$ 是平展的，就可以在更熟悉的框架中描述这一分类。若 $F$ 的特征为零，则总是如此；若 $|A|$ 与特征互素，也是如此，但并非总是如此。这一恼人困难出现的最简单情形，是特征 $2$ 的域上的群 $G=\operatorname{GL}(2)$。此时 $A$ 是由 $z^2=1$ 定义的 $\operatorname{GL}(1)$ 的子群。为简化写作，我们假定 $A$ 在 $F$ 上平展；然而要建立完整的理论，必须去掉这一条件。

只有当自守谱含有离散部分时，我们才能有效地使用迹公式。若 $G$ 的中心含有一个 $F$-分裂环面，则并非如此；此时商 $G(F)\backslash G(\mathbb A_F)$ 的体积无限。于是必须或者像 Arthur 那样把 $G(\mathbb A_F)$ 换成一个更小的群，或者把 $G(F)$ 换成一个更大的群。归根结底，两种选择的差别很小。我们更愿意把 $G(F)$ 换成更大的群，这在我们看来是最优雅的解决办法。

记 $Z_{\mathrm{sp}}$ 为 $Z$ 中最大的 $F$-分裂环面，$Z'_{\mathrm{sp}}$ 为 $G$ 的最大的、是 $F$-分裂环面的商。由此得到的同态 $Z_{\mathrm{sp}}\to Z'_{\mathrm{sp}}$ 是一个同源，记其核为 $Z''_{\mathrm{sp}}$。记 $s_G$ 为 $Z_{\mathrm{sp}}$ 与 $Z'_{\mathrm{sp}}$ 的秩。选取子群
\[
Z_{s_G}\subset Z_{\mathrm{sp}}(\mathbb A_F)
\]
使得商 $Z_{s_G}Z_{\mathrm{sp}}(F)\backslash Z_{\mathrm{sp}}(\mathbb A_F)$ 是紧的。群 $Z_{s_G}$ 与 $Z''_{\mathrm{sp}}$ 的交必然平凡，因此 $Z_{s_G}\to Z'_{\mathrm{sp}}(\mathbb A_F)$ 是把 $Z_{s_G}$ 嵌入 $Z'_{\mathrm{sp}}(\mathbb A_F)$。对 $G$ 的任意子群 $H$，我们记
\[
H_+(F)=Z_{s_G}H(F).
\]
特别地，$G_+(F)=Z_{s_G}G(F)$。商
\[
Z_{s_G}Z(F)\backslash Z(\mathbb A_F)=Z_+(F)\backslash Z(\mathbb A_F)
\]
是紧群。

不可约自守表示 $\pi$ 的中心特征定义了 $Z(F)\backslash Z_+(F)=Z_{s_G}$ 的一个特征 $\chi$。存在 $Z'_{\mathrm{sp}}(F)\backslash Z'_{\mathrm{sp}}(\mathbb A_F)$ 的特征 $\chi'$，其在 $Z_{s_G}$ 上的限制为 $\chi$。把 $\pi$ 换成 $\pi\otimes\chi'^{-1}$，就得到中心特征在 $Z_+(F)$ 上平凡的自守表示。因此我们可以只限于具有这一性质的自守表示，并且只考虑商
\[
G_+(F)\backslash G(\mathbb A_F),
\]
它对 $G(\mathbb A_F)$ 上的不变测度有有限体积。

\subsection{测度}
\srcpage{19}

众所周知，如何把测度与实微分流形上的体积形式相联系。只要在非 Archimedes 局部域上选定一个起实数域上 Lebesgue 测度作用的测度，同样的做法对局部域上的流形也适用；对阿代尔环也是如此。这些测度可以按一种典范的方式定义，理解这一点对我们非常重要，但即使在最常被引用的文献（[W1, W2]）中，它也没有得到很好的解释。由于对我们而言，有一份精确而又尽可能简短的参考材料确实重要，我们在此重述这些定义。于是我们先在 阿代尔环 $\mathbb A_F$ 与所有位 $v\in|F|$ 的局部域 $F_v$ 上固定相容的测度。

局部紧群 $\mathbb A_F$ 上的不变 Haar 测度在相差一个常数意义下是确定的。由于 $\mathbb A_F$ 对离散群 $F$ 的商是紧群，可以按唯一的方式归一化测度，使商 $F\backslash\mathbb A_F$ 的测度为 $1$。本文其余部分就采用这个不变测度 $dx$；但最好不要直接用"商的测度为 $1$"这一条件来定义它，而应从局部域 $F_v$（$v$ 为 $F$ 的位）上的测度 $dx_v$ 出发，而这些测度本身又由一个连续整体加法特征定义。正是这一系列定义及其逻辑在 [W2] 中没有得到很好的解释，并在记忆层面、最终在逻辑层面上导致困难，当我们到达约化群上的迹公式时尤其如此。我们从一个整体特征 $\chi$ 开始，它给出每个位 $v$ 的特征 $\chi_v$；由 $\chi_v$ 定义测度 $dx_v$。若 $|F|$ 表示所有这些位的集合，则 $dx=\otimes_{v\in|F|}dx_v$ 将是 $\mathbb A_F$ 上的测度。

考虑连续特征 $\chi:\mathbb A_F\to\mathbb C^\times$ 的集合，满足：对一切 $b\in F$ 有 $\chi(b)=1$，且对一切 $b\in F$ 有 $\chi(bb')=1$ 当且仅当 $b'\in F$。这是群 $F^\times$ 下的一个主齐性空间。令人不快的是，这样的特征的存在性在函数域情形与数域情形下的验证方式不同；对那些钟爱罗塞塔石碑的人来说，这甚至可能令人困惑。

\srcpage{20}
在两种情形下，任何特征 $\chi$ 都由其局部分量定义：$\chi(a)=\prod_v\chi_v(a_v)$。

若 $F$ 是有限域 $\kappa$（$q$ 个元素）上射影光滑曲线 $X$ 的有理函数域，则选取 $X$ 上一个非零亚纯微分 $\omega$。于是有特征 $\mathbb A_F\to\kappa$，由
\[
x\longmapsto\operatorname{res}(x\omega)=\sum_v\operatorname{tr}_{F_v/\kappa}\bigl(\operatorname{res}_v(x\omega)\bigr)
\]
定义，且当 $x=b\in F$ 时它平凡。由此得到特征 $\chi:\mathbb A_F\to\mathbb C^\times$，定义为
\begin{equation}\tag{3.2}
x\longmapsto\exp\left(\frac{2\pi i}{p}\operatorname{tr}_{\kappa/\mathbb F_p}\bigl(\operatorname{res}(x\omega)\bigr)\right).
\end{equation}
若 $F=\mathbb Q$ 且 $x\in\mathbb A_F$，令
\[
\chi_0(x)=\prod_v\chi_v(x_v),
\]
其中 $\chi_\infty(x_\infty)=\exp(-2\pi ix_\infty)$，而 $\chi_p(x_p)=\exp(2\pi ix')$，只要 $x'\in\mathbb Q$ 的分母是 $p$ 的幂且满足 $|x'-x_p|_p\le1$。最后，对 $\mathbb Q$ 的任意有限扩张 $F$，令
\begin{equation}\tag{3.3}
\chi(x)=\chi_0\bigl(\operatorname{tr}_{F/\mathbb Q}(x)\bigr).
\end{equation}

固定了 $\mathbb A_F$ 的特征 $\chi$ 之后，对每个位 $v$ 就得到 $F_v$ 的特征 $\chi=\chi_v$。于是 $F_v$ 上存在唯一的 Haar 测度 $dx=dx_v$，它关于 Fourier 变换
\[
\hat f(y)=\int f(x)\chi(xy)\,dx
\]
是自对偶的，也就是说，关于这个测度有 $\hat{\hat f}(x)=f(-x)$。若 $v$ 是非 Archimedes 位，且
\[
d_v^{-1}=d_v^{-1}(\chi)=\{x\mid\chi(xy)=1\ \ \forall y\in\mathcal O_v\},
\]
则自对偶测度给 $\mathcal O_v$ 的测度为 $N d_v^{-1/2}$。理想 $d_v$ 与局部不同有关，但并不等于这个不同；它与特征相关联。由于它的存在，没有典范的局部计算。不过 $d_v$ 几乎处处等于 $\mathcal O_v$，因此几乎处处有典范的公式。通常的计算表明：在实位，与特征 $\chi(x)=\exp(2\pi iyx)$ 相伴的自对偶测度是 $|y|^{1/2}dx$；在复位，与 $\chi(z)=\exp(2\pi i\operatorname{Re}(\bar wz))$ 相伴的测度是 $|w|dxdy$，其中 $z=x+iy$，$w=u+iv$，$|w|=\sqrt{u^2+v^2}$。

\srcpage{21}
固定了整体特征 $\chi$ 之后，我们同时固定了局部自对偶测度 $dx_v$ 以及 $\mathbb A_F$ 上的整体测度 $dx=\prod_v dx_v$，后者也是自对偶的。由乘积公式 $\prod_v|b|_v=1$（$b\in F^\times$），后一测度与 $\chi$ 的选取无关。由于 $F\backslash\mathbb A_F$ 的对偶群是 $F$，由紧群上 Fourier 变换的一般理论容易推出
\begin{equation}\tag{3.4}
\meas(F\backslash\mathbb A_F)=1
\end{equation}
以及
\[
\sum_{b\in F}\hat f(-b)=\sum_{b\in F}f(b).
\]
这就是 Poisson 公式，它给出与 $F$ 相伴的 $L$ 函数的函数方程。

设 $X$ 是 $F$ 上维数为 $n$ 的光滑代数簇。$X$ 上一个处处非零的 $n$-次微分形式 $\omega$ 定义了拓扑空间 $X(\mathbb A_F)$ 上的一个测度 $|\omega|$。对每一点 $x\in X(F_v)$，$X(F_v)$ 中存在 $x$ 的一个解析邻域，同构于以 $a_1,\dots,a_n$ 为坐标的开多圆盘。在这个多圆盘上，$n$-形式写作 $\omega=f\,da_1\wedge\cdots\wedge da_n$，其中 $f$ 是多圆盘上可逆的解析函数。测度 $|f|\,da_1\cdots da_n$ 转移到 $x$ 的邻域上，实际上与局部坐标的选取无关。这个测度只依赖于 $n$-形式 $\omega$，故记作 $|\omega|_v$。若再给出 $X$ 在 $\mathcal O_v$ 上的一个光滑模型 $\mathcal X$，并假定 $n$-形式 $\omega$ 延拓为概形 $\mathcal X$ 上一个处处非零的不变 $n$-形式，则有公式
\begin{equation}\tag{3.5}
\int_{X(F_v)}\mathbf 1_{X(\mathcal O_v)}|\omega|_v
=q_v^{-n}N d_v^{-n/2}|\mathcal X(\kappa_v)|,
\end{equation}
其中 $\kappa_v$ 是 $\mathcal O_v$ 的剩余域，$q_v$ 是它的基数，$|\mathcal X(\kappa_v)|$ 表示 $\mathcal X$ 的 $\kappa_v$-点个数，最后 $\mathbf 1_{X(\mathcal O_v)}$ 表示紧集 $\mathcal X(\mathcal O_v)$ 的特征函数。这些原理的最简单情形，显然就是有限维向量空间上的不变形式。

\srcpage{22}
设 $G$ 是代数群。\footnote{尽管这些预备定义对任意群都成立，我们一开始就假定它是约化的。它不必是拟分裂的；但到达稳定迹公式时，我们加上这一条件。我们提醒，稳定公式的理论，也就是内窥，把任意约化群上的自守形式研究归结为拟分裂群的理论。然而，我们并不假定它们在 $F$ 的非分歧扩张上分裂。

在几何学家中，回避分歧或只处理温和分歧的做法非常普遍。从几何观点看，只考虑 Hecke 算子的理论、回避非不变调和分析的许多问题，或者回避在算术中如此普遍的分歧，当然有好处。不过，这些问题照亮了许多结构性问题，当我们在考察迹公式时，它们对我们很重要。} $G$ 上一个不变 $n$-形式的给出，等价于在 $F$-向量空间 $\wedge^n\mathfrak g$（维数为一，$\mathfrak g$ 是 $G$ 的 Lie 代数）中给出一个非零向量 $\omega$。与此相伴的不变 $n$-形式在 $F_v$ 上处处非零，在 $\mathcal O_v$ 上几乎处处非零。对每个位 $v$，它定义 $G(F_v)$ 上的不变测度 $|\omega|_v$。一般不能在 阿代尔群上取测度 $\prod_{v\in|F|}|\omega|_v$，因为公式 (3.5) 蕴含无穷乘积 $\prod_v\mu_v$ 一般发散，其中 $\mu_v$ 表示紧子群 $G(\mathcal O_v)$ 的测度：
\[
\mu_v=\int_{G(F_v)}\mathbf 1_{G(\mathcal O_v)}|\omega|_v=\int_{G(\mathcal O_v)}|\omega_v|.
\]
特别地，若 $G$ 是 $\mathbb Q$ 上的乘法群，则对每个素数 $p$，局部测度 $\mu_p=1-p^{-1}$，它是 Riemann $\zeta$ 函数的 $p$ 因子之逆在 $1$ 处的值。

为了在取乘积之前修正局部测度、从而得到定义良好的测度，需要使用多种局部测度。尽管并非绝对必要，有些作者不仅引入整体乘积测度，还引入它的重归一化。本文使用一种记号来区分所有这些测度，并为后续文章中将要用到的两三个测度给出更简单的记号。首先令 $d^{\mathrm{geom}}g_v=|\omega|_v$，它是由形式 $\omega$ 直接给出的测度，显然依赖于这个形式。设 $\sigma_G$ 是定义在代数闭包 $\bar F$ 上的特征 $G\to\mathbb G_m$ 的 Galois 模构成的表示；它们也定义在可分闭包上。与 $\sigma_G$ 上的表示相伴的 Artin $L$ 函数

\srcpage{23}
在 $\operatorname{Re}s>1$ 上有 Euler 乘积展开
\[
L(s,\sigma_G)=\prod_vL_v(s,\sigma_G).
\]
$G(F_v)$ 上的归一化局部测度定义为
\begin{equation}\tag{3.6}
dg_v=d^{\mathrm{norm}}g_v=L_v(1,\sigma_G)\,d^{\mathrm{geom}}g_v.
\end{equation}
一旦基本定义被充分理解，我们对此后一个测度就只用记号 $dg_v$，不再加额外的下标，因为它就是主要的局部测度。

乘积测度
\begin{equation}\tag{3.7}
dg=d^{\mathrm{prod}}g=\bigotimes_{v\in|F|}dg_v
\end{equation}
由于乘积公式 $\prod_{v\in|F|}|c|_v=1$（对一切 $c\in F^\times$）而与 $\omega$ 无关。根据 Ono（见 [O3] 附录 2），这个乘积测度是定义良好的。若 $f$ 是形如 $f=\otimes_{v\in|F|}f_v$ 的函数，其中 $f_v$ 光滑紧支且几乎处处等于 $G(\mathcal O_v)$ 的特征函数，则无穷乘积
\[
\prod_{v\in|F|}\int_{G(F_v)}f_v\,dg_v
\]
绝对收敛。

人们常常再加一个整体重归一化，但我们宁愿不这样做。不过我们还是解释一下。Galois 模 $\sigma_G$ 的平凡部分的秩等于 $Z_{\mathrm{sp}}$ 的秩 $s_G$，且 $L$ 函数 $L(s,\sigma_G)$ 在 $s=1$ 有 $s_G$ 阶极点。若
\begin{equation}\tag{3.8}
\rho_G=\lim_{s\downarrow1}(s-1)^{s_G}L(s,\sigma_G),
\end{equation}
则
\begin{equation}\tag{3.9}
d\tilde g=\rho_G^{-1}d^{\mathrm{prod}}g
\end{equation}
是双重归一化的测度，但我们不常使用它。我们仍然引入两个测度：乘积测度，记作 "measprod"，由 $d^{\mathrm{prod}}g$ 定义且局部测度已归一化；以及简单地记作 "meas" 的双重归一化测度。这不是因为它更重要，而是因为它更流行。这可能会使本文中前一个不那么方便的记号频繁出现，但实际上我们会避免它。乘积测度的主要优点，是它在 阿代尔框架中更易于操作，而这个框架是 $G$ 上调和分析的传统框架，我们也将把它引入 Steinberg–Hitchin 基。

\srcpage{24}
我们选取离散子群 $G_+(F)$，使得商 $G_+(F)\backslash G(\mathbb A_F)$ 有有限测度。它的测度
\[
\meas\bigl(G_+(F)\backslash G(\mathbb A_F)\bigr)
\]
出现在迹公式中，容易由 Ono 的文章（[O3]）中定义并计算的 Tamagawa 测度 $\tau(G)$ 算出。引入正合列
\[
\{1\}\to G_1\to G\to G_2\to\{1\},
\]
其中 $G_1$ 连通，$G_2$ 是秩为 $s_G$ 的分裂环面。$G_2(\mathbb A_F)$ 在 $G(\mathbb A_F)$ 中的原像就是 [O3] 中的群 $G^1_{\mathbb A}$。测度
\[
\meas\bigl(G_+(F)\backslash G(\mathbb A_F)\bigr),
\]
关于归一化整体测度计算，等于 $G(F)\backslash G^1_{\mathbb A}$ 的测度与 $G^1_{\mathbb A}\backslash G_{\mathbb A}$ 对 $G_+(F)$ 的像（一个离散群）的商的测度之积。这个商 $Q$ 是紧的，在数域情形同构于 $Z_{s_G}\backslash\mathbb R^{s_G}$，在函数域情形同构于有限群。这些测度按 [O3] 的定义计算。因此
\begin{equation}\tag{3.10}
\meas\bigl(G_+(F)\backslash G(\mathbb A_F)\bigr)=\tau(G)\int_Qdt.
\end{equation}
最后这个不起眼的因子 $m_G=\int_Qdt$ 依赖于 $G_+(F)$ 的选取，从而依赖于 $Z_{s_G}$。因子 $\tau(G)$ 是 Tamagawa 数。\footnote{在函数域情形，$Z_{s_G}$ 的几何意义以及 $m_G$ 的意义，我们留给读者去解释。}

若群 $G$ 是环面 $T$，则由 (3.5)，无穷乘积
\begin{equation}\tag{3.11}
\prod_v\int_{T(F_v)}f_v|\omega_v|
\end{equation}
的因子几乎处处等于 $1/L_v(1,\sigma_T)$。若 $T$ 是各向异性的，即 $s_T=0$，则 Dirichlet 定理及其推广蕴含乘积
\begin{equation}\tag{3.12}
\prod_v\frac{1}{L_v(1,\sigma_T)}
\end{equation}
在按 $q_v$（$q_v=|F_v|$）的大小排列 $v$ 时条件收敛，且此时 (3.12) 等于 $1/\rho_G$。因此 (3.11) 给出的结果与 $\int f\,d\tilde g$ 只相差有限个局部因子之积。
